% Actividad de Clase 2 - Identificación del Objeto de Prueba (SIDEM)
% Transcripción de la hoja entregada en clase, con F06 y el punto 6 completados.
% Punto 7 con las cuatro casillas marcadas y los nombres de los integrantes.
\documentclass[11pt,letterpaper]{article}
\usepackage[utf8]{inputenc}
\usepackage[T1]{fontenc}
\usepackage{lmodern}
\usepackage{amssymb}
\usepackage[margin=2cm,headheight=14pt]{geometry}
\usepackage{array}
\usepackage{tabularx}
\usepackage{booktabs}
\usepackage{fancyhdr}
\usepackage{titlesec}
\usepackage{hyperref}

\hypersetup{hidelinks,pdftitle={Actividad de Clase 2 - SIDEM},pdfauthor={Equipo SIDEM}}

\pagestyle{fancy}
\fancyhf{}
\fancyhead[L]{\small Plan de Pruebas de Software -- Actividad de Clase 2}
\fancyhead[R]{\small ISO/IEC/IEEE 29119-3}
\fancyfoot[C]{\small Página \thepage}
\renewcommand{\headrulewidth}{0.5pt}

\titleformat{\section}{\large\bfseries}{\thesection.}{0.5em}{}
\titlespacing*{\section}{0pt}{14pt}{6pt}

\newcolumntype{Y}{>{\raggedright\arraybackslash}X}
\newcolumntype{P}[1]{>{\raggedright\arraybackslash}p{#1}}
\newcommand{\si}{$\boxtimes$\,Sí\enspace$\square$\,No}
\newcommand{\vacia}{$\square$\,Sí\enspace$\square$\,No}
\newcommand{\linea}[1]{\underline{\makebox[6.5cm][l]{#1}}}

\begin{document}

\begin{center}
{\Large\bfseries ACTIVIDAD DE CLASE 2}\par\vspace{3pt}
{\large Identificación del Objeto de Prueba (\textit{Test Items})}\par\vspace{2pt}
{\small Norma ISO/IEC/IEEE 29119-3:2013 \quad|\quad Duración: 60 minutos}
\end{center}
\vspace{6pt}

\noindent
\begin{tabularx}{\textwidth}{|Y|Y|}
\hline
\textbf{EQUIPO DE TRABAJO} (Desarrolladores) & \textbf{EQUIPO REVISOR} (Evaluadores) \\
\hline
Sistema: \textbf{SIDEM} \newline Fecha: \textbf{24-09-2026} & Equipo: \textbf{RESTful Booker} \\
\hline
Integrantes: \newline
1. Nicolás Castaño Durango \newline
2. Julián Ramírez Molina \newline
3. Sebastián Díaz Jiménez \newline
4. Marcos Duque Jaramillo
&
Integrantes revisores: \newline
1. Susana Romero García \newline
2. Santiago Montoya Restrepo \newline
3. Miguel Ángel Vallejo Giraldo \newline
4. Luisa Fernanda Tamayo Arias \\
\hline
\end{tabularx}

\section{Propósito y Marco Teórico}
En la Actividad 1 se definieron la metodología, el objetivo y el alcance general del plan de pruebas. En esta actividad se avanza a \textbf{identificar los elementos concretos del sistema que serán objeto de prueba} (\textit{test items}).

De acuerdo con la norma \textbf{ISO/IEC/IEEE 29119-3:2013}, los \textit{test items} corresponden a los módulos, componentes o funcionalidades susceptibles de ser evaluados. En esta fase no se diseñan casos de prueba detallados: se delimita qué se probará, bajo qué condiciones y con qué prioridad.

\section{Identificación y Delimitación del Objeto de Prueba}
Retomen el sistema de la Actividad 1 e identifiquen entre \textbf{6 y 7 funcionalidades clave}. Indiquen qué hace cada elemento, quién interactúa con él, si se incluye en el alcance y la razón técnica.

\vspace{4pt}
\noindent
{\small
\begin{tabularx}{\textwidth}{@{}P{0.8cm}P{2.6cm}YP{2.3cm}P{2.0cm}Y@{}}
\toprule
\textbf{ID} & \textbf{Elemento / Funcionalidad} & \textbf{¿Qué hace?} & \textbf{Actor principal} & \textbf{¿En alcance?} & \textbf{Justificación de selección o exclusión} \\
\midrule
F01 & Inicio de sesión & Permite el ingreso al sistema con usuario y contraseña. & Todos los usuarios & \si & Controla el acceso general al sistema. \\
\addlinespace
F02 & Gestión de usuarios y roles & Crea usuarios, asigna roles y define qué puede hacer cada uno. & Administrador & \si & Garantiza que cada usuario acceda a lo que le corresponde. \\
\addlinespace
F03 & Registro del seguimiento & Permite evaluar los criterios de cada curso en cada revisión. & Revisor & \si & Es la función principal del sistema y la fuente de los datos de métricas. \\
\addlinespace
F04 & Semáforo de cumplimiento & Clasifica en un color según el porcentaje de cumplimiento. & Líderes & \si & Resume el estado del seguimiento y apoya la toma de decisiones. \\
\addlinespace
F05 & Panel de métricas & Muestra estadísticas del seguimiento por revisión. & Líderes / administrador & \si & Permite validar que la información sea confiable. \\
\addlinespace
F06 & Exportación de informes & Genera un informe en Excel con los resultados del seguimiento. & Revisor / líderes & \si & Es la salida formal de la información y debe coincidir con lo registrado en el sistema. \\
\addlinespace
F07 & & & & \vacia & \\
\bottomrule
\end{tabularx}
}

\newpage
\section{Condiciones de Prueba: Precondiciones y Criterios de Aceptación}
Seleccionen \textbf{tres funcionalidades incluidas} en el alcance y describan las condiciones iniciales de entrada y el resultado esperado que validaría su correcto funcionamiento.

\vspace{4pt}
\noindent
\begin{tabularx}{\textwidth}{@{}P{0.9cm}YY@{}}
\toprule
\textbf{ID} & \textbf{Precondiciones / Datos de entrada} & \textbf{Resultado esperado / Criterio de aceptación} \\
\midrule
F02 & Existe un usuario con un rol de solo consulta. Ese usuario intenta modificar un registro de seguimiento. & El sistema impide la modificación y muestra un mensaje de acceso denegado. El usuario solo puede ver la información. \\
\addlinespace
F03 & El revisor ha iniciado sesión y existe un curso asignado en una revisión abierta. El revisor marca los criterios y escribe una observación. & La información queda guardada, se asocia al curso y a la revisión correctos y se ve igual al volver a consultarla. \\
\addlinespace
F04 & Existen tres cursos evaluados en la misma revisión: uno con el 80\,\% de los criterios cumplidos, otro con el 50\,\% y otro con el 40\,\%. & El semáforo muestra verde, amarillo y rojo, respectivamente. \\
\bottomrule
\end{tabularx}

\section{Priorización del Objeto de Prueba}
Identifiquen \textbf{tres elementos prioritarios} para el proceso de pruebas. Justifiquen su elección según criterios técnicos o de negocio: \textit{Criticidad para el usuario, Complejidad técnica/integración, Frecuencia de uso o Impacto ante fallas.}

\vspace{4pt}
\noindent
\begin{tabularx}{\textwidth}{@{}P{1.9cm}P{4.2cm}Y@{}}
\toprule
\textbf{Prioridad} & \textbf{Elemento (ID y nombre)} & \textbf{Justificación del criterio de priorización} \\
\midrule
1 (Alta) & F02 Gestión de usuarios y roles & Tiene el mayor impacto ante fallas: un error en los permisos permitiría ver o modificar información de otros usuarios o departamentos. \\
\addlinespace
2 (Media) & F03 Registro del seguimiento & Es la funcionalidad de mayor frecuencia de uso, y los errores en el registro afectan los resultados posteriores. \\
\addlinespace
3 (Media) & F04 Semáforo de cumplimiento & Es crítico para el usuario, porque los líderes toman decisiones con base en este resultado. Un cálculo incorrecto genera conclusiones equivocadas. \\
\bottomrule
\end{tabularx}

\section{Revisión Cruzada entre Equipos (Coevaluación de Pares)}
Intercambien el documento con otro equipo. El objetivo es verificar si una persona ajena al desarrollo comprende qué partes del sistema serán objeto de prueba y por qué.

\vspace{4pt}
\noindent
\begin{tabularx}{\textwidth}{@{}Ycc@{}}
\toprule
\textbf{Criterio de evaluación} & \textbf{Sí} & \textbf{No} \\
\midrule
Los elementos identificados representan funcionalidades reales y concretas del sistema. & $\boxtimes$ & $\square$ \\
Los actores principales de cada funcionalidad están claramente definidos. & $\boxtimes$ & $\square$ \\
La delimitación del alcance es coherente con la Actividad 1 y sus exclusiones están justificadas. & $\boxtimes$ & $\square$ \\
Las precondiciones y criterios de aceptación reflejan la lógica esperada del software. & $\boxtimes$ & $\square$ \\
La priorización se justifica técnicamente (criticidad, complejidad o riesgo de negocio). & $\boxtimes$ & $\square$ \\
\bottomrule
\end{tabularx}

\vspace{8pt}
\noindent\textbf{Dos fortalezas o aspectos adecuadamente definidos:}\par
1. Gestión de usuarios y roles.\par
2. Semáforo de cumplimiento.

\vspace{6pt}
\noindent\textbf{Dos oportunidades de mejora o aspectos a precisar:}\par
1. No se encontraron oportunidades o aspectos a mejorar.\par
2. No se encontraron oportunidades o aspectos a mejorar.

\section{Ajustes del Equipo y Reflexión Metacognitiva}
\noindent\textbf{Ajustes realizados tras la coevaluación:}\par
El equipo revisor no reportó oportunidades de mejora. En la autorrevisión se completó F06 (Exportación de informes) para cumplir el mínimo de 6 funcionalidades; los demás elementos se mantienen sin cambios.

\vspace{6pt}
\noindent\textbf{Reflexión pedagógica:} ¿Por qué es un error metodológico diseñar casos de prueba detallados antes de delimitar y priorizar el objeto de prueba?\par
Porque se invierte esfuerzo sin criterio: se prueban funciones fuera del alcance y las más críticas quedan con poca cobertura. La norma ISO/IEC/IEEE 29119-3 pide definir primero qué se probará y con qué prioridad.

\section{Registro de Entrega y Firmas}
\noindent
\begin{tabular}{@{}p{8cm}p{8cm}@{}}
$\boxtimes$ Objeto de prueba delimitado. & $\boxtimes$ Coevaluación realizada con equipo par. \\
$\boxtimes$ Condiciones de prueba descritas. & $\boxtimes$ Ajustes e integrantes registrados. \\
\end{tabular}

\vspace{14pt}
\noindent\textbf{Firmas de Integrantes del Equipo:}\par\vspace{22pt}
\noindent 1. \linea{Marcos Duque Jaramillo}\hfill 3. \linea{Nicolás Castaño Durango}\par\vspace{22pt}
\noindent 2. \linea{Sebastián Díaz Jiménez}\hfill 4. \linea{Julián Ramírez Molina}

\end{document}
